% Part: sets-functions-relations
% Chapter: relations
% Section: equivalence-relations
% Hindi translation (hi-Deva-IN) of Open Logic commit 9620cc73f9c8e0ad003c514a5d3748f29611c4c0.

\documentclass[../../../include/open-logic-section]{subfiles}

\begin{document}

\olfileid[hi]{sfr}{rel}{eqv}

\olsection{तुल्यता संबंध}

किसी समुच्चय पर तत्समक संबंध स्वतुल्य, सममित और संक्रामक होता है। जिन
संबंधों~$R$ में ये तीनों गुणधर्म होते हैं, वे बहुत सामान्य हैं।

\begin{defn}[तुल्यता संबंध]
जो संबंध $R \subseteq A^2$ स्वतुल्य, सममित और संक्रामक हो, उसे
\emph{तुल्यता संबंध} कहते हैं। $x$ और $y$, जो~$A$ के !!^{element}s हैं,
\emph{$R$-तुल्य} कहलाते हैं यदि~$Rxy$।
\end{defn}

तुल्यता संबंधों से \emph{तुल्यता वर्ग} की धारणा उत्पन्न होती है। कोई
तुल्यता संबंध अपने प्रांत को अलग-अलग खंडों में विभाजित कर देता है। प्रत्येक
खंड के भीतर सभी वस्तुएँ एक-दूसरे से संबंधित होती हैं; और अलग-अलग खंडों की
कोई भी वस्तुएँ आपस में संबंधित नहीं होतीं। कभी-कभी इन खंडों की बात
\emph{सीधे} करना ही उपयोगी होता है। इस उद्देश्य से हम निम्न परिभाषा देते
हैं:

\begin{defn}\ollabel{def:equivalenceclass}
मान लें कि $R \subseteq A^2$ एक तुल्यता संबंध है। प्रत्येक $x \in A$ के
लिए, $x$ का~$A$ के भीतर \emph{तुल्यता वर्ग} समुच्चय $\equivrep{x}{R}
= \Setabs{y \in A}{Rxy}$ है। $A$ का~$R$ के अधीन \emph{विभाग समुच्चय}
$\equivclass{A}{R} = \Setabs{\equivrep{x}{R}}{x \in A}$ है, अर्थात इन
तुल्यता वर्गों का समुच्चय।
\end{defn}

अगला परिणाम यह सिद्ध करके तुल्यता वर्ग की परिभाषा का औचित्य स्थापित करता
है कि तुल्यता वर्ग वास्तव में~$A$ के विभाजन के खंड ही हैं:

\begin{prop}
यदि $R \subseteq A^2$ एक तुल्यता संबंध है, तो $Rxy$ तभी और केवल तभी जब
$\equivrep{x}{R} = \equivrep{y}{R}$।
\end{prop}

\begin{proof}
बाएँ से दाएँ दिशा के लिए मान लें कि $Rxy$, और $z \in
\equivrep{x}{R}$ हो। तब परिभाषा से $Rxz$। चूँकि $R$ एक तुल्यता संबंध है,
$Ryz$। (इसे विस्तार से देखें: चूँकि $Rxy$ और~$R$ सममित है, इसलिए $Ryx$;
और चूँकि $Rxz$ और~$R$ संक्रामक है, इसलिए~$Ryz$।) अतः $z \in \equivrep{y}{R}$।
सामान्यीकरण करने पर,
$\equivrep{x}{R} \subseteq \equivrep{y}{R}$। ठीक इसी प्रकार,
$\equivrep{y}{R} \subseteq \equivrep{x}{R}$। इसलिए विस्तारिता से
$\equivrep{x}{R} =
\equivrep{y}{R}$।

दाएँ से बाएँ दिशा के लिए मान लें कि $\equivrep{x}{R} =
\equivrep{y}{R}$। चूँकि $R$ स्वतुल्य है, $Ryy$, इसलिए $y \in
\equivrep{y}{R}$। इसलिए $y \in \equivrep{x}{R}$ भी है, क्योंकि हमारी
धारणा है कि $\equivrep{x}{R} = \equivrep{y}{R}$। अतः $Rxy$।
\end{proof}

\begin{ex}
तुल्यता संबंधों का एक अच्छा उदाहरण मापांक अंकगणित से मिलता है। किन्हीं भी
$a$, $b$ और $n \in \PosInt$ के लिए, $a \equiv_n b$ तभी और केवल तभी मानें
जब $a$ को~$n$ से भाग देने पर वही शेषफल मिले जो $b$ को~$n$ से भाग देने पर
मिलता है। (कुछ अधिक सांकेतिक रूप में: $a \equiv_n b$ तभी और केवल तभी जब
किसी $k \in
\Int$ के लिए $a - b = kn$।) अब $\equiv_n$ किसी भी~$n$ के लिए एक तुल्यता
संबंध है। और ठीक $n$ भिन्न तुल्यता वर्ग~$\equiv_n$ से उत्पन्न होते हैं;
अर्थात $\equivclass{\Nat}{\equiv_n}$ में $n$ !!{element}s हैं। ये हैं:
$n$ से निःशेष विभाज्य संख्याओं का समुच्चय, अर्थात
$\equivrep{0}{\equiv_n}$; $n$ से भाग देने पर शेषफल~$1$ देने वाली संख्याओं
का समुच्चय, अर्थात $\equivrep{1}{\equiv_n}$; \ldots; और $n$ से भाग देने
पर शेषफल~$n-1$ देने वाली संख्याओं का समुच्चय, अर्थात
~$\equivrep{n-1}{\equiv_n}$।
\end{ex}

\begin{prob}
सिद्ध कीजिए कि $\equiv_n$ किसी भी $n \in
\PosInt$ के लिए एक तुल्यता संबंध है, और
$\equivclass{\Nat}{\equiv_n}$ में ठीक $n$ अवयव हैं।
\end{prob}

\end{document}
